\begin{lemma}[Characterization of the chosen energy-normalized gradient]
Let $v \in L^2(\gamma_n)$ and let $\Gamma \in L^2(\gamma_n;\mathbb{R}^n)$, and
suppose $\sqrt2\,\Gamma$ is a weak Euclidean gradient of $v$ in the sense of
\noderef{HasGaussWeakGrad}.  Then $v \in W^{1,2}(\gamma_n)$ in the sense of
\noderef{SobolevW12}, and $Dv = \Gamma$ with $D$ the operator of \noderef{GaussD}.
The conjunction is the whole content: by \noderef{GaussD} the value $Dv$ is set to
$0$ by convention when $v \notin W^{1,2}(\gamma_n)$, so the equality $Dv = \Gamma$
on its own asserts nothing about $\Gamma$ unless the membership is also asserted.
\end{lemma}
\begin{proof}
The two conclusions are proved from different inputs: the membership comes from the
existential form of \noderef{SobolevW12}, which the hypothesis witnesses directly,
while the identification $Dv = \Gamma$ requires \noderef{GaussWeakGradUnique}, which
is what upgrades ``$\Gamma$ is \emph{a} field with the defining property'' to
``$\Gamma$ is \emph{the} field with it''.

\emph{Step 1: membership.}
By \noderef{SobolevW12}, $v \in W^{1,2}(\gamma_n)$ means precisely that
$v \in L^2(\gamma_n)$ and that there exists some
$\Gamma' \in L^2(\gamma_n;\mathbb{R}^n)$ for which $\sqrt2\,\Gamma'$ is a weak
Euclidean gradient of $v$ in the sense of \noderef{HasGaussWeakGrad}.  The
hypotheses of the present lemma supply $v \in L^2(\gamma_n)$ and exhibit
$\Gamma' = \Gamma$ as such a witness.  Hence $v \in W^{1,2}(\gamma_n)$.

\emph{Step 2: the selected gradient is itself a weak gradient.}
By \noderef{GaussD}, for $v \in W^{1,2}(\gamma_n)$ the element $Dv$ of
$L^2(\gamma_n;\mathbb{R}^n)$ is characterized by the property that $\sqrt2\,(Dv)$
is a weak Euclidean gradient of $v$ in the sense of \noderef{HasGaussWeakGrad}.
Since Step 1 has established $v \in W^{1,2}(\gamma_n)$, this property is available:
\begin{equation*}
  \sqrt2\,(Dv) \ \text{is a weak Euclidean gradient of } v. \tag{1}
\end{equation*}
In reading $(1)$ and the hypothesis we are free to replace the pointwise scalar
multiple $x \mapsto \sqrt2\,(Dv)(x)$ by the element $\sqrt2\,(Dv)$ of
$L^2(\gamma_n;\mathbb{R}^n)$, and likewise for $\Gamma$: the two agree
$\gamma_n$-almost everywhere, and by \noderef{GaussWeakGradAECongr} the
weak-gradient relation is unchanged when either argument is replaced by a
$\gamma_n$-almost-everywhere equal function.

\emph{Step 3: uniqueness.}
Both $\sqrt2\,(Dv)$ and $\sqrt2\,\Gamma$ lie in $L^2(\gamma_n;\mathbb{R}^n)$,
being scalar multiples of elements of that space, and by $(1)$ and the hypothesis
both are weak Euclidean gradients of the same function $v$.  Therefore
\noderef{GaussWeakGradUnique}, applied to the function $v$ and to the two fields
$\sqrt2\,(Dv)$ and $\sqrt2\,\Gamma$, gives
\[
  \sqrt2\,(Dv) = \sqrt2\,\Gamma
  \qquad\text{in } L^2(\gamma_n;\mathbb{R}^n).
\]
Since $\sqrt2 \neq 0$ and $L^2(\gamma_n;\mathbb{R}^n)$ is a real vector space,
scalar multiplication by $\sqrt2$ is injective on it, so $Dv = \Gamma$.
\end{proof}
